\setchapterpreamble[u]{\margintoc}
\chapter{Avaliações}
\labch{cet615-avaliacoes}

A disciplina teve duas avaliações individuais, entregues pelo Moodle como
arquivo TXT ou PDF; a nota final é a média das duas, e a média para
aprovação é 7. A professora pediu respostas assertivas: só os comandos
pedidos, sem dados de teste e sem alterar o esquema físico fornecido.

\section{Primeira avaliação}

Entregue até 14/11/2026, às 9h. As duas questões usam o esquema de um
sistema de venda de obras literárias usadas, o mesmo do mapeamento do
\refch{cet615-conceitos}:

\begin{itemize}
	\item \texttt{cliente(login, senha, nome, endereco,
	qtd\_obras\_vendidas, qtd\_obras\_compradas)};
	\item \texttt{obras\_literarias(codobra, nome, editora, preco, tipo,
	login\_vendedor, data)}, com \texttt{tipo} 0 para livro e 1 para
	periódico, e \texttt{login\_vendedor} referenciando \texttt{cliente};
	\item \texttt{livro(codobra, ISBN)} e \texttt{periodico(codobra,
	numero, volume)}, cada um com \texttt{codobra} referenciando a obra;
	\item \texttt{compra(codcompra, data, login\_comprador, codobra)}, com
	um índice único em \texttt{codobra}: cada obra é comprada no máximo
	uma vez.
\end{itemize}

\begin{figure*}[h]
\centering
\begin{tikzpicture}[
	tabela/.style={rectangle split, rectangle split parts=2, draw=bdroxo, thick,
		rectangle split part fill={bdroxo, white}, font=\scriptsize, align=left,
		inner sep=2pt, text width=31mm},
	fk/.style={-{Stealth[length=2mm]}, draw=bdazul, thick}]
	\node[tabela] (cli) at (0,0) {\textcolor{white}{\textbf{cliente}}
		\nodepart{two} \underline{login}\\ senha, nome, endereco\\ qtd\_obras\_vendidas\\ qtd\_obras\_compradas};
	\node[tabela] (obr) at (5.2,0) {\textcolor{white}{\textbf{obras\_literarias}}
		\nodepart{two} \underline{codobra}\\ nome, editora, preco\\ tipo\\ login\_vendedor (FK)\\ data};
	\node[tabela] (com) at (10.4,0) {\textcolor{white}{\textbf{compra}}
		\nodepart{two} \underline{codcompra}\\ data\\ login\_comprador (FK)\\ codobra (FK, única)};
	\node[tabela] (liv) at (3.0,-3.2) {\textcolor{white}{\textbf{livro}}
		\nodepart{two} \underline{codobra} (FK)\\ ISBN};
	\node[tabela] (per) at (7.4,-3.2) {\textcolor{white}{\textbf{periodico}}
		\nodepart{two} \underline{codobra} (FK)\\ numero, volume};
	\draw[fk] (obr.west) -- node[rotulo, above] {vende} (cli.east);
	\draw[fk] (com.west) -- node[rotulo, above] {obra} (obr.east);
	\draw[fk] (com.north) -- ++(0,0.4) -| node[rotulo, above, pos=0.25] {comprador} (cli.north);
	\draw[fk] (liv.north) -- (liv.north |- obr.south);
	\draw[fk] (per.north) -- (per.north |- obr.south);
\end{tikzpicture}
\caption{Tabelas do esquema da primeira avaliação, desenhadas a partir do
esquema lógico. As setas vão da chave estrangeira para a tabela
referenciada; chaves primárias sublinhadas.}
\end{figure*}

\subsection{Questão 1 (4,0): visão}

Pedido: uma visão que, para cada periódico já comprado, mostre o login e
o nome do vendedor, o código e o nome da obra, a data em que foi posta à
venda e a data da compra; e um exemplo de consulta à visão.

\begin{lstlisting}
create view periodicos_comprados
as
select c.login login_vendedor, c.nome nome_vendedor,
       o.codobra, o.nome nome_obra,
       o.data data_venda, cp.data data_compra
from cliente c inner join obras_literarias o
         on c.login = o.login_vendedor
         inner join periodico p
         on o.codobra = p.codobra
         inner join compra cp
         on o.codobra = cp.codobra
go

select * from periodicos_comprados
where login_vendedor = 'ana'
order by data_compra
\end{lstlisting}

Decisões:

\begin{itemize}
	\item A junção com \texttt{periodico} restringe a obra aos periódicos;
	\texttt{o.tipo = 1} daria o mesmo resultado com dados consistentes. A
	junção com \texttt{compra} restringe às obras já compradas.
	\item As duas colunas \texttt{nome} e as duas \texttt{data} recebem
	apelidos, porque teriam o mesmo nome na visão.
	\item A visão traz só o que foi pedido; o filtro e a ordenação ficam na
	consulta de exemplo, como recomendado no \refch{cet615-visoes}.
\end{itemize}

\subsection{Questão 2 (6,0): procedimento armazenado}

Pedido: um procedimento para a transação de cadastro de livros, que
retorne 1 em caso de sucesso e 0 caso contrário, com um exemplo de
execução que mostre o retorno. Os vendedores já estão cadastrados.

\begin{lstlisting}
create procedure ins_livro
@codobra int,
@nome char(40),
@editora char(40),
@preco numeric(6,2),
@login_vendedor char(8),
@data date,
@isbn numeric(13,0)
as
begin transaction
insert into obras_literarias
values (@codobra, @nome, @editora, @preco, 0, @login_vendedor, @data)
if @@rowcount > 0
begin
      insert into livro
      values (@codobra, @isbn)
      if @@rowcount > 0
      begin
            commit transaction
            return 1
      end
      else
      begin
            rollback transaction
            return 0
      end
end
else
begin
      rollback transaction
      return 0
end
go

declare @ret int
exec @ret = ins_livro 10, 'Dom Casmurro', 'Garnier', 25.00, 'ana',
            '2026-11-07', 9788535910663
print @ret
\end{lstlisting}

Decisões, seguindo o \refch{cet615-procedimentos}:

\begin{itemize}
	\item A transação tem as duas operações do cadastro: a obra genérica e
	o livro específico, nessa ordem por causa da chave estrangeira.
	\item O código da obra é passado uma vez só, e o \texttt{tipo} é fixo
	(0), não parâmetro.
	\item Os tipos dos parâmetros são os mesmos das colunas do esquema
	físico.
	\item A contagem \texttt{qtd\_obras\_vendidas} do vendedor ficou de
	fora. Na apresentação da avaliação, a professora pediu a transação
	mais compacta possível e anunciou que a segunda avaliação usaria o
	mesmo esquema. Totais derivados como esse são o exemplo clássico de
	gatilho (\refch{cet615-gatilhos}).
\end{itemize}

\subsection{Verificação}

As duas respostas foram testadas no SQL Server 2022 com dados de exemplo
antes da entrega. A visão devolveu só os periódicos comprados, deixando
de fora um periódico não comprado e um livro comprado. O procedimento
retornou 1 num cadastro válido. Também retornou 0 com código de obra
repetido e com vendedor inexistente, sem deixar nada gravado e sem
transação aberta. Nos casos de falha, o SQL Server mostra a mensagem do
erro (violação de chave primária ou estrangeira) antes do retorno 0. É o
comportamento esperado do padrão com \texttt{@@rowcount}.

\section{Segunda avaliação}

\textit{Será incluída depois da entrega, prevista para 21/11/2026, às 9h.}
